\documentclass[11pt]{article}
\usepackage[margin=0.85in]{geometry}
\usepackage[T1]{fontenc}
\usepackage{lmodern}
\usepackage{microtype}
\usepackage{amsmath,amssymb,mathtools}
\usepackage{booktabs,array,tabularx,longtable}
\usepackage{xcolor}
\usepackage{hyperref}
\usepackage{enumitem}
\usepackage{fancyhdr}
\usepackage{titlesec}
\usepackage{lastpage}
\hypersetup{colorlinks=true,urlcolor=blue!55!black,linkcolor=blue!55!black,citecolor=blue!55!black}
\setlist[itemize]{leftmargin=*,itemsep=2pt,topsep=3pt}
\setlist[enumerate]{leftmargin=*,itemsep=3pt,topsep=3pt}
\renewcommand{\arraystretch}{1.18}
\setlength{\parskip}{5pt}
\setlength{\parindent}{0pt}
\titleformat{\section}{\large\bfseries\color{blue!35!black}}{\thesection.}{0.55em}{}
\titleformat{\subsection}{\normalsize\bfseries}{\thesubsection}{0.55em}{}
\pagestyle{fancy}
\fancyhf{}
\fancyhead[L]{\small MATH 440 | Group 5}
\fancyhead[R]{\small NBA Gravity Mini-project 2}
\fancyfoot[C]{\small Page \thepage\ of \pageref{LastPage}}
\newcommand{\todo}[1]{\textcolor{red!65!black}{\textbf{[TO CONFIRM:} #1\textbf{]}}}
\newcommand{\draft}[1]{\textcolor{blue!50!black}{\textbf{[PROPOSED:} #1\textbf{]}}}
\newcommand{\nba}{\mathrm{NBA}}

\begin{document}
\begin{center}
    {\LARGE\bfseries Investigating NBA Offensive Gravity}\\[5pt]
    {\large Mini-project 2 --- Shared Mathematical and Research Plan}\\[8pt]
    \textbf{MATH 440 | Group 5 | Josue and Emma}\\[3pt]
    \href{https://github.com/katnissali/NBA-Gravity}{GitHub repository}\\[3pt]
    \textit{Working draft v0.8 --- October 22, 2026}\\[3pt]
    \textbf{Next presentation: Thursday, October 22, 2026, at 2:00 p.m.}
\end{center}

\noindent\colorbox{blue!5}{\parbox{\dimexpr\linewidth-2\fboxsep\relax}{\textbf{Purpose and current status.} This working report documents the research question, six measured predictor variables, source provenance, mathematical definitions, verified aggregation and dataset integration, and preliminary OLS evaluation. The final player-level file \texttt{player\_gravity\_model\_data.csv} contains 234 official GRAV players, including 175 with all six predictors. Results include one held-out split and five-fold cross-validation on the training set; substantial fold-to-fold variation warrants cautious interpretation. Blue labels denote proposals; red labels denote items requiring confirmation.}}

\setcounter{tocdepth}{1}
\tableofcontents
\medskip

\section{Research question, definitions, and propositions}
\textbf{Research question:} Can publicly available NBA statistics be combined into
an interpretable mathematical model that estimates player offensive gravity,
and how closely do the predictions and ranking agree with the NBA's official
2025--26 regular-season Player Gravity metric?

\textbf{Offensive gravity:} Defensive attention that a player attracts beyond
what would ordinarily be expected from location and playing context. The NBA's
tracking-based GRAV metric measures defensive pressure above expectation
\cite{nba-gravity-intro,nba-gravity-board}. Our model predicts that
metric from publicly accessible performance statistics; it does not reproduce
the NBA's tracking algorithm.

\textbf{Working hypothesis (H):} Stronger gravity-related offensive indicators
are associated with higher official NBA GRAV values. These propositions concern
\emph{associations}, not demonstrated causes.

\textbf{Propositions to test (in the project's official factor order):}
\begin{enumerate}[label=\textbf{P\arabic*:}]
\item Players with higher team offensive rating ON/OFF differentials tend to
have higher official GRAV scores.
\item Players with higher \emph{teammates-only} effective field-goal percentage
ON/OFF differentials tend to have higher official GRAV scores.
\item Players averaging more assists per game tend to have higher official
GRAV scores.
\item Players with higher opponent timeout-response rates per 36 playing
minutes tend to have higher official GRAV scores.
\item Players scoring more points per 36 minutes in qualifying clutch
situations tend to have higher official GRAV scores.
\item Players with greater \emph{Offensive Win Shares (OWS)} tend to have
higher official GRAV scores.
\item \textbf{TBD (future factor):} A positive association between
teammates' open-shot percentage ON/OFF differential and official GRAV will be
considered when an operational definition and reliable tracking data are available.
\item A model combining several suitably defined indicators gives a player
ranking positively associated with the NBA Gravity ranking.
\end{enumerate}
Propositions P1--P6 correspond to six collected candidate predictors. P7
is a deferred proposal, not a testable claim with the currently assembled
data. P8 retains the original combined-model ranking-agreement proposition.
Better prediction than a simple baseline is an additional \emph{evaluation
criterion}, not a demonstrated finding.

\section{Collected data: 2025--26 regular season}
All six predictor datasets and the official NBA GRAV target concern the
2025--26 regular season. The individual source files have now been
\emph{combined into one player-level dataset}. An initial OLS fit,
held-out evaluation, and five-fold cross-validation on the training set
have been completed. Fold-to-fold variability is reported below.

\begin{center}
\small
\begin{tabularx}{\linewidth}{@{}p{0.39\linewidth}rX@{}}
\toprule
\textbf{Dataset} & \textbf{Rows} & \textbf{Purpose / verified coverage}\\
\midrule
\texttt{player\_stats.csv} & 582 & NBA per-game statistics; the \texttt{AST} column is assists per game.\\
\texttt{team\_stats.csv} & 30 & Original season-level team data, retained for context.\\
\texttt{player\_on\_off\_stats.csv} & 1,362 & Earlier NBA on/off dataset; not used as either final ON/OFF predictor.\\
\texttt{player\_off\_rating\_on\_off.csv} & 701 & Player--team observations; 661 valid ON/OFF pairs.\\
\texttt{player\_teammate\_efg\_on\_off.csv} & 701 & Teammates-only ON/OFF eFG; 661 valid player--team pairs.\\
\texttt{player\_timeout\_rates.csv} & 582 & Timeout-response rates using all 1,230 regular-season games.\\
\texttt{player\_clutch\_rates.csv} & 246 & Qualifying players with at least 30 total clutch minutes.\\
\texttt{player\_win\_shares\_with\_nba\_ids.csv} & 582 & Basketball-Reference Offensive Win Shares, matched to NBA IDs.\\
\texttt{player\_gravity\_official.csv} & 234 & Distinct official NBA player IDs with nonmissing \texttt{GRAV}.\\
\midrule
\texttt{player\_on\_off\_aggregated.csv} & 602 & One row per player after ON-minutes weighting of factors 1 and 2.\\
\texttt{player\_gravity\_model\_data.csv} & \textbf{234} & \textbf{Final merged GRAV-plus-six-factor dataset.}\\
\bottomrule
\end{tabularx}
\end{center}

The two final ON/OFF input datasets each have \textbf{701 player--team rows
for 602 distinct players}. Among them, \textbf{87 players} have multiple
recorded team stints. The weighted aggregation produced 602 distinct player
records, with valid factor~1 and factor~2 values for 582 players apiece.
Among the 234 GRAV players, \textbf{24 have multiple recorded team stints}.
The 234 GRAV players all have both aggregated ON/OFF factors.

The NBA's \texttt{GravityLeaders} endpoint supplies the official target:
\texttt{AVGGRAVITYSCORE}, preserved as \texttt{GRAV}. The different field
\texttt{GRAVITYSCORE} is \emph{not} used as the target. The assembled
player-level table uses \texttt{PLAYER\_ID} as its unique join key.
No missing factor has been replaced with zero.

\textbf{Verified complete-case coverage:} factors 1, 2, 3, 4, and 6 are
available for \textbf{234/234} GRAV players; factor 5 (clutch scoring) is
available for \textbf{175/234}. Thus the merged file contains 234 rows,
\textbf{175 complete six-factor observations} and 59 with a missing clutch
measurement. This is a data-availability finding, not a regression result.

\section{Official factor order and candidate variables}
The six active predictors are indexed $j=1,\dots,6$ in the order agreed by
the group. The proposed seventh factor, teammates' open-shot percentage, is
\textbf{TBD} and excluded from the first regression.

\small
\begin{longtable}{@{}p{0.25\linewidth}p{0.44\linewidth}p{0.24\linewidth}@{}}
\toprule
\textbf{Factor} & \textbf{Candidate statistical variable} & \textbf{Collection status}\\
\midrule
\endfirsthead
\toprule
\textbf{Factor} & \textbf{Candidate statistical variable} & \textbf{Collection status}\\
\midrule
\endhead
1. Team offensive rating ON/OFF & $X_{i1}=\mathrm{OffRtg}_{i,\mathrm{ON}}-\mathrm{OffRtg}_{i,\mathrm{OFF}}$ (points per 100 possessions). & Aggregated; 234/234 GRAV players.\\
\addlinespace
2. Teammates-only eFG\% ON/OFF & $X_{i2}=100(eFG_{i,\mathrm{ON}}^{\mathrm{tm}}-eFG_{i,\mathrm{OFF}}^{\mathrm{tm}})$ (percentage points). & Aggregated; 234/234 GRAV players.\\
\addlinespace
3. Assists & $X_{i3}=\mathrm{AST}_{i}^{\mathrm{per\ game}}$ (assists per game; confirmed representation). & Merged; 234/234 GRAV players.\\
\addlinespace
4. Opponent timeout responses & $X_{i4}=36N_i/M_i$ (responses per 36 season minutes). & Merged; 234/234 GRAV players.\\
\addlinespace
5. Clutch scoring & $X_{i5}=36\,\mathrm{PTS}^{\mathrm{clutch}}_i/\mathrm{MIN}^{\mathrm{clutch}}_i$ (clutch points per 36 minutes). & Merged; 175/234 GRAV players.\\
\addlinespace
6. Offensive Win Shares & $X_{i6}=\mathrm{OWS}_i$ (season Offensive Win Shares). & Merged; 234/234 GRAV players.\\
\addlinespace
7. Teammate open-shot percentage ON/OFF & Potential future ON/OFF differential; definition of ``open-shot percentage'' still TBD. & Deferred; excluded from v1.\\
\bottomrule
\end{longtable}
\normalsize

\textbf{Terminology:} The fifth collected variable measures \emph{clutch
scoring rate}, not ``clutch share'' or a percentage of team clutch points.
Changing it to a share would require a different data definition and approval.

\textbf{Interpretation:} All six predictors are proposed proxies for gravity.
Offensive rating, teammate eFG, and timeout responses may be influenced by
lineups, opposing defenses, and game context. Assists, clutch scoring, and Offensive Win Shares may measure broader offensive skill or playing time rather than
purely defensive attention. None alone establishes causality.

\section{Mathematical definitions}
Let $i$ index a player in the final analysis sample, and let $j\in\{1,\dots,6\}$
index the six proposed predictors in Section~3. The symbol $X_{ij}$ denotes
the \emph{numerical value of predictor $j$ for player $i$}, before
standardization. An ON or OFF quantity refers to the player's own team while
that player is on or off the court. Unless otherwise specified, counts and
rates refer to the 2025--26 regular season. For the two ON/OFF predictors, the formulas below first describe a single
team stint. The verified aggregation rule later in this section combines
these stints into one player-season measurement.

\subsection{Factor 1: Team offensive rating ON/OFF}
\textbf{Offensive rating (OffRtg)} is the number of team points scored per
100 offensive possessions. For player $i$ and court-status
$s\in\{\mathrm{ON},\mathrm{OFF}\}$, define
\begin{equation}
\mathrm{OffRtg}_{i,s}
=100\frac{\mathrm{PTS}^{\mathrm{team}}_{i,s}}
{\mathrm{POSS}^{\mathrm{team}}_{i,s}},
\qquad \mathrm{POSS}^{\mathrm{team}}_{i,s}>0.
\label{eq:offrtg}
\end{equation}
Here $\mathrm{PTS}^{\mathrm{team}}_{i,s}$ is the number of points scored by
the player's team during condition $s$, and
$\mathrm{POSS}^{\mathrm{team}}_{i,s}$ is its offensive-possession count in
that condition. We use the NBA's reported \texttt{OFF\_RATING} fields from
\texttt{TeamPlayerOnOffSummary} (\texttt{Advanced}, \texttt{Totals});
we do not estimate those possession counts from rounded per-game data.
The \textbf{first predictor} is
\begin{equation}
\boxed{X_{i1}=\mathrm{OffRtg}_{i,\mathrm{ON}}
-\mathrm{OffRtg}_{i,\mathrm{OFF}}.}
\label{eq:offrtg-diff}
\end{equation}
Its unit is \emph{team points per 100 possessions}. A positive difference
means the observed team OffRtg is higher with player $i$ present; it does
not establish the player's causal contribution.

\subsection{Factor 2: Teammates-only eFG\% ON/OFF}
\textbf{Effective field-goal percentage (eFG)} adjusts field-goal accuracy
for the extra point value of a made three-pointer:
\begin{equation}
 eFG=\frac{\mathrm{FGM}+0.5\,\mathrm{FG3M}}{\mathrm{FGA}},
 \qquad \mathrm{FGA}>0.
 \label{eq:efg}
\end{equation}
Here \texttt{FGM} means made field goals, \texttt{FG3M} means made
three-point field goals, and \texttt{FGA} means attempted field goals.
The resulting eFG is a proportion: 0.55 represents 55\%.

For player $i$, let $\mathrm{FGM}^{\mathrm{team}}_{i,s}$,
$\mathrm{FG3M}^{\mathrm{team}}_{i,s}$, and
$\mathrm{FGA}^{\mathrm{team}}_{i,s}$ denote the team's shooting totals
in condition $s$. Let $\mathrm{FGM}^{\mathrm{self}}_i$,
$\mathrm{FG3M}^{\mathrm{self}}_i$, and $\mathrm{FGA}^{\mathrm{self}}_i$
denote the player's own shot totals \emph{in that team stint}. Teammates'
ON-court efficiency excludes those personal shots:
\begin{equation}
\begin{split}
 eFG^{\mathrm{tm}}_{i,\mathrm{ON}}=
 \frac{
 (\mathrm{FGM}^{\mathrm{team}}_{i,\mathrm{ON}}-\mathrm{FGM}^{\mathrm{self}}_i)
 +0.5(\mathrm{FG3M}^{\mathrm{team}}_{i,\mathrm{ON}}-\mathrm{FG3M}^{\mathrm{self}}_i)
 }{
 \mathrm{FGA}^{\mathrm{team}}_{i,\mathrm{ON}}-\mathrm{FGA}^{\mathrm{self}}_i
 }.
\end{split}
\label{eq:efg-on}
\end{equation}
The OFF-court team shots already exclude player $i$:
\begin{equation}
 eFG^{\mathrm{tm}}_{i,\mathrm{OFF}}=
 \frac{\mathrm{FGM}^{\mathrm{team}}_{i,\mathrm{OFF}}
 +0.5\,\mathrm{FG3M}^{\mathrm{team}}_{i,\mathrm{OFF}}}
 {\mathrm{FGA}^{\mathrm{team}}_{i,\mathrm{OFF}}}.
 \label{eq:efg-off}
\end{equation}
Both denominators must be positive. Team shooting counts come from
\texttt{TeamPlayerOnOffDetails}; personal shooting totals come from
\texttt{TeamPlayerDashboard}. The \textbf{second predictor} is
\begin{equation}
 \boxed{X_{i2}=100\left(eFG^{\mathrm{tm}}_{i,\mathrm{ON}}
 -eFG^{\mathrm{tm}}_{i,\mathrm{OFF}}\right).}
 \label{eq:efg-diff}
\end{equation}
$X_{i2}$ is measured in \emph{percentage points}. For example, a difference
of 0.04 between the underlying eFG proportions corresponds to 4 percentage
points, not a 4\% relative increase. The calculation excludes the focal
player's shooting contribution from the ON value.

\subsection{Factor 3: Assists}
An \textbf{assist} is a credited pass leading directly to a teammate's made
basket under NBA recording rules. The source
\texttt{player\_stats.csv} contains \texttt{AST} as a \emph{per-game}
statistic, as confirmed during data preparation. The \textbf{third
predictor} is
\begin{equation}
 \boxed{X_{i3}=\mathrm{AST}^{\mathrm{per\ game}}_i.}
 \label{eq:assists}
\end{equation}
Thus $X_{i3}$ is the player's season-average assists per game, with unit
\emph{assists/game}. This representation is convenient but may favor
players with more playing time per game. The merged file preserves the
source's per-game representation under the column \texttt{AST\_PER\_GAME};
we do not divide it by games a second time or convert it to per 36 minutes.

\subsection{Factor 4: Opponent timeout responses}
We define one \textbf{timeout response} for player $i$ when an opposing
team's timeout is the \emph{next recorded play-by-play event} after that
player's made field goal. A timeout by the shooter's team, or a timeout
after an intervening event, does not qualify. Define $N_i$ as the season
count of qualifying responses and $M_i$ as total regular-season minutes
played by player $i$. The \textbf{fourth predictor} is
\begin{equation}
 \boxed{X_{i4}=36\frac{N_i}{M_i}},\qquad M_i>0.
 \label{eq:timeouts}
\end{equation}
Its unit is \emph{opponent timeout responses per 36 playing minutes}.
The project analyzed all 1,230 regular-season games using
\texttt{PlayByPlayV3} and season-total player minutes. The definition
records event order, not evidence that the player's basket caused the
timeout.

\subsection{Factor 5: Individual clutch scoring rate}
The NBA's standard \textbf{clutch} condition is the final five minutes
of the fourth quarter or overtime when the score difference is within
five points. For player $i$, let $\mathrm{PTS}^{\mathrm{clutch}}_i$
be accumulated points and $\mathrm{MIN}^{\mathrm{clutch}}_i$
be accumulated playing minutes in those situations. The
\textbf{fifth predictor} is
\begin{equation}
 \boxed{X_{i5}=36\frac{\mathrm{PTS}^{\mathrm{clutch}}_i}
 {\mathrm{MIN}^{\mathrm{clutch}}_i}},
 \qquad \mathrm{MIN}^{\mathrm{clutch}}_i\ge 30.
 \label{eq:clutch}
\end{equation}
The unit is \emph{clutch points per 36 clutch minutes}. The 30-minute minimum reduces extremely small-sample rates. The
saved qualifying clutch file contains 246 player rows. Of the 234 official
GRAV players, 175 have qualifying clutch data and 59 do not. Missing
clutch values are retained as missing, not interpreted as zero points.
Selecting players with 30 or more clutch minutes may also introduce
sample-selection bias, because such players have had more opportunities
to appear in close games.

\subsection{Factor 6: Offensive Win Shares}
\textbf{Offensive Win Shares (OWS)} is a published Basketball-Reference
estimate of the number of team wins attributable to a player's offensive
contribution over the season. It is \emph{not} a raw count of team wins,
team wins divided by roster size, or total Win Shares.
Basketball-Reference defines
\begin{equation}
 \mathrm{WS}_i=\mathrm{OWS}_i+\mathrm{DWS}_i,
 \label{eq:winshares}
\end{equation}
where $\mathrm{WS}_i$ denotes total Win Shares and $\mathrm{DWS}_i$
denotes Defensive Win Shares. Our \textbf{sixth predictor} uses only the
offensive component:
\begin{equation}
 \boxed{X_{i6}=\mathrm{OWS}_i.}
 \label{eq:ws-predictor}
\end{equation}

\textbf{Meaning of the published calculation.} For a modern NBA season,
Basketball-Reference describes the following methodology
\cite{winshares}:
\begin{align}
 \mathrm{MO}_i & = \mathrm{PP}_i - 0.92\,\mathrm{PPP}_{L}\,\mathrm{OP}_i,
 \label{eq:ows-marginal}\\
 \mathrm{MPW}_{i} & =0.32\,\mathrm{PPG}_{L}
 \left(\frac{\mathrm{Pace}_{\mathrm{team}(i)}}{\mathrm{Pace}_{L}}\right),
 \label{eq:ows-points}\\
 \mathrm{OWS}_i & = \frac{\mathrm{MO}_i}{\mathrm{MPW}_i}.
 \label{eq:ows-method}
\end{align}
Here $\mathrm{PP}_i$ is the player's estimated \emph{points produced},
$\mathrm{OP}_i$ is the estimated number of offensive possessions attributable
to the player, and $\mathrm{PPP}_{L}$ is league-wide points per possession.
$\mathrm{PPG}_{L}$ is league-wide points per game;
$\mathrm{Pace}_{\mathrm{team}(i)}$ and $\mathrm{Pace}_{L}$ are team and
league pace. $\mathrm{MO}_i$ is \emph{marginal offense}, and
$\mathrm{MPW}_i$ is \emph{marginal points per win}. These equations explain
the source metric; \textbf{our group did not independently reconstruct
points produced, possessions, or OWS}. We extracted the published \texttt{OWS}
values and used them as supplied. Negative OWS values are possible; they
must not be converted to zero.

\textbf{Collection and matching.} The 2025--26 values were taken from
Basketball-Reference's advanced player statistics for the season
\cite{br-advanced} and preserved in
\texttt{player\_win\_shares\_with\_nba\_ids.csv}. The CSV has
\textbf{582 distinct NBA player IDs} and includes both the
Basketball-Reference player key (\texttt{BR\_PLAYER\_ID}) and the NBA
player key (\texttt{PLAYER\_ID}). It retains player names, team labels,
games, minutes and name/team matching indicators to make the identity
mapping auditable. In the saved file, 562 names matched exactly; 9 required
punctuation/accent normalization and 11 are recorded as verified name
variants. The team-check field records 506 same-team matches, 72
multi-team Basketball-Reference records, and 4 different team labels.
These are recorded checks, not proof that every original match was
manually verified. All \textbf{234 official GRAV players} have a matching,
nonmissing OWS value. The final dataset stores the season OWS value as
\texttt{OWS} without converting it to OWS per game or WS/48.

OWS is a \emph{cumulative} estimate; higher values can reflect both stronger
offensive contribution and more playing time. It is a proposed correlate
of gravity rather than a direct tracking measurement of defensive
attention or a demonstrated causal factor.

\subsection{Factor 7: Teammate open-shot percentage (TBD)}
This proposed seventh factor has not been collected. ``Open-shot
percentage'' could mean the \emph{fraction of shots attempted with a
sufficiently distant defender}, or \emph{shooting accuracy on those
attempts}. These are distinct statistics. We will not assign it a
numeric variable or include it in the six-predictor model until its
meaning and the necessary ON/OFF teammate tracking data are established.

\subsection{Official NBA target variable: GRAV}
Let $Y_i$ denote player $i$'s \textbf{official Average Gravity score}
for the 2025--26 NBA regular season. In the NBA's
\texttt{GravityLeaders} data, the leaderboard's numeric \texttt{GRAV}
value is \texttt{AVGGRAVITYSCORE}, and we record
\begin{equation}
 \boxed{Y_i=\mathrm{GRAV}_i=\mathrm{AVGGRAVITYSCORE}_i.}
 \label{eq:grav-target}
\end{equation}
The official league data contains 234 players, 234 distinct NBA player
IDs, and no missing values in this field. The reference scores range
from approximately $-10.225$ to $16.196$ in the extracted data. The
NBA describes its normalized gravity metric on a theoretical $-100$ to
$100$ scale \cite{nba-gravity-intro,nba-gravity-board,nba-gravity-api}.
Each $Y_i$ has now been joined to our candidate predictors via
\texttt{PLAYER\_ID}. Preliminary OLS coefficients have been estimated on
the training portion of the complete-case sample; held-out results are
reported in Section~\ref{sec:model-evaluation}.


\subsection{Combining ON/OFF measurements across multiple teams}
The official GRAV dataset has one observation per player, whereas the
ON/OFF source files have one observation per player--team combination.
Let $i$ identify a player, let $t$ identify a team, and let
$\mathcal{T}_i$ be the set of teams recorded for player $i$.
For each stint, define the two recorded differences
\begin{align}
 d_{i1t}&=\mathrm{OffRtg}_{i,t,\mathrm{ON}}
           -\mathrm{OffRtg}_{i,t,\mathrm{OFF}},\label{eq:stint-off}\\
 d_{i2t}&=100\bigl(eFG^{\mathrm{tm}}_{i,t,\mathrm{ON}}
           -eFG^{\mathrm{tm}}_{i,t,\mathrm{OFF}}\bigr).
           \label{eq:stint-efg}
\end{align}
The first difference has units of points per 100 possessions; the second
has units of percentage points. Define $m_{it}^{(j)}$ as player $i$'s
ON-court minutes for team $t$ in predictor $j$'s source file, for
$j\in\{1,2\}$. Let $\mathcal{V}_{ij}\subseteq\mathcal{T}_i$
be the teams with a recorded $d_{ijt}$ and a positive, finite value
$m_{it}^{(j)}$.

We use the following \emph{minutes-weighted mean of stint differences}
for the season-level first and second predictors:
\begin{equation}
 \boxed{
 X_{ij}=\frac{\displaystyle\sum_{t\in\mathcal{V}_{ij}}
 m_{it}^{(j)}d_{ijt}}{
 \displaystyle\sum_{t\in\mathcal{V}_{ij}}m_{it}^{(j)}}
 },\qquad j\in\{1,2\}.
 \label{eq:multi-team-aggregation}
\end{equation}
If $\mathcal{V}_{ij}$ is empty, $X_{ij}$ is missing rather than zero.
A team record with no ON minutes or no valid ON/OFF difference receives
no weight. A single valid team stint retains its original difference.
The same calculation is carried out separately for $j=1$ and $j=2$;
thus the eligible team sets may differ if one source has a missing value.

This rule was implemented by our aggregation script,
\texttt{aggregate\_gravity\_on\_off.py}.
Its output, \texttt{player\_on\_off\_aggregated.csv}, has
602 unique players; 582 have valid aggregated values
for each of the two factors, including \textbf{all 234} official GRAV players.

A diagnostic also compared the NBA's reported ON minutes in the two source
files. Of 661 player--team pairs with both minute values, 646 differ by
more than 0.01 minutes; the median absolute difference was 0.2433 minutes,
the maximum 0.5 minutes, and the median relative difference 0.0321\%.
These small differences are \emph{consistent with} distinct reporting
precision (for example, whole minutes versus fractional minutes), rather
than substantive differences in exposure. Each predictor uses minutes
from its own source file, as specified in Equation~\eqref{eq:multi-team-aggregation}.

This rule is a reproducible player-level \emph{summary}, not an exact
season-long on/off statistic. In particular, averaging stint differences
is not identical to recomputing overall eFG from pooled shot totals or
overall offensive rating from pooled points and possessions. The source
OFF-court figures may also include games outside the player's tenure
with a particular team. Minutes weighting alone cannot remove these
contextual limitations. The raw player--team records were retained.
Alternative aggregation methods remain a possible sensitivity analysis.
Of 234 official GRAV players, 24 have multiple recorded team stints.
No regression is implied by this aggregation.

\section{Final merged player-level dataset and quality checks}
The integration was executed using
\texttt{merge\_gravity\_model\_data.py}. Starting from the 234-row official
GRAV file, the script performed player-ID one-to-one left joins with
the already-aggregated ON/OFF file, per-game player stats, timeout rates,
qualifying clutch rates, and OWS data. It did not merge the original
701-row player--team files directly, fill missing values, or fit a model.
The saved result is \textbf{\texttt{data/processed/player\_gravity\_model\_data.csv}}.

\begin{center}
\small
\begin{tabularx}{\linewidth}{@{}p{0.39\linewidth}p{0.31\linewidth}X@{}}
\toprule
\textbf{Column} & \textbf{Role / units} & \textbf{GRAV coverage}\\
\midrule
\texttt{PLAYER\_ID}, \texttt{PLAYER\_NAME} & Identifying variables & 234/234\\
\texttt{GRAV} & Response $Y_i$ (official NBA Gravity) & 234/234\\
\texttt{OFF\_RATING\_DIFF} & $X_{i1}$; points per 100 possessions & 234/234\\
\texttt{EFG\_DIFF\_PP} & $X_{i2}$; percentage points & 234/234\\
\texttt{AST\_PER\_GAME} & $X_{i3}$; assists/game & 234/234\\
\texttt{TIMEOUTS\_PER\_36} & $X_{i4}$; responses/36 season minutes & 234/234\\
\texttt{CLUTCH\_PTS\_PER\_36} & $X_{i5}$; points/36 clutch minutes & 175/234\\
\texttt{OWS} & $X_{i6}$; offensive win-share estimate & 234/234\\
\texttt{ALL\_SIX\_AVAILABLE} & Completeness indicator (Boolean) & 175 \texttt{True}\\
\bottomrule
\end{tabularx}
\end{center}

For each player $i$, the completeness flag is
\begin{equation}
 C_i=\begin{cases}
 1,&\text{if }X_{ij}\text{ is observed for every }j=1,\ldots,6,\\
 0,&\text{otherwise}.
 \end{cases}
 \label{eq:complete-case}
\end{equation}
The CSV column \texttt{ALL\_SIX\_AVAILABLE} stores this flag as
\texttt{True} or \texttt{False}. The script verified unique player IDs and
retained all 234 target rows. Its final console validation reported
\textbf{175 players with six observed predictors} and \textbf{59 missing
at least one factor}. In this dataset the missing factor for those 59
players is clutch scoring. The complete-case sample is
$\mathcal C=\{i:C_i=1\}$, with $|\mathcal C|=175$. The complete-case share
is $175/234\approx74.8\%$.

\textbf{Important interpretation:} a missing clutch measurement means
the player does not meet our qualifying data rule; it does \emph{not}
mean the player's clutch scoring rate is zero. The 175-player regression
sample would exclude these 59 players, so test results may not generalize
to all 234 players on the NBA GRAV leaderboard.

\textbf{Sharing and reproducibility:} The model CSV exists locally in
\texttt{data/processed/}. For this report, we assume all group members have
access to the shared GitHub repository and the required project files.
This access assumption was provided by the group and was not independently
tested here. The collection checkpoints are not needed to fit the
already-assembled player-level CSV. If any scripts are rerun after file
reorganization, their input/output paths should be checked first.

\section{Constructing an interpretable gravity model}
\textbf{Primary method:} Multiple linear regression with six predictors.
\textbf{Benchmark:} An equally weighted index of those same standardized predictors.
The first model fit and held-out evaluation are complete and reported
below alongside five-fold cross-validation results. The substantial
fold-to-fold variability means these results do not establish stable
generalization.

\subsection{Standardization}
The six indicators have different units. For player $i$ and predictor $j$,
\begin{equation}
Z_{ij}=\frac{X_{ij}-\mu_j^{\mathrm{train}}}{s_j^{\mathrm{train}}},
\qquad s_j^{\mathrm{train}}>0.
\label{eq:zscore}
\end{equation}
Means and standard deviations must be computed \emph{only on training players}
and then applied unchanged to held-out players. Any orientation changes, missing
value rules, and zero-variance exclusions must be declared in advance.

\subsection{Equal-weight benchmark}
Given an agreed orientation for every predictor, define the index
\begin{equation}
Q_i=\frac{1}{6}\sum_{j=1}^{6}Z_{ij}.
\label{eq:equal}
\end{equation}
$Q_i$ is a dimensionless score. It can be compared with NBA ranks using Spearman
correlation. For a fair MAE comparison on the NBA GRAV scale, a training-only
calibration such as $\widehat Y_i^{\mathrm{equal}}=a+bQ_i$ would be needed.

\subsection{Multiple linear regression}
Our candidate model is
\begin{equation}
Y_i=\beta_0+\sum_{j=1}^{6}\beta_j Z_{ij}+\varepsilon_i,
\qquad
\widehat Y_i=\widehat\beta_0+\sum_{j=1}^{6}\widehat\beta_j Z_{ij}.
\label{eq:ols}
\end{equation}
For training-player index set $\mathcal T$, ordinary least squares minimizes
\begin{equation}
\underset{\beta_0,\boldsymbol\beta}{\operatorname{arg\,min}}
\sum_{i\in\mathcal T}
\left(Y_i-\beta_0-\sum_{j=1}^{6}\beta_jZ_{ij}\right)^2.
\label{eq:least-squares}
\end{equation}
For a full-column-rank training design matrix $X$ including the intercept,
the normal equations are
\begin{equation}
X^{\mathsf T}X\widehat{\boldsymbol\theta}=X^{\mathsf T}\mathbf Y,
\qquad
\widehat{\boldsymbol\theta}=(X^{\mathsf T}X)^{-1}X^{\mathsf T}\mathbf Y,
\label{eq:normal}
\end{equation}
where $\boldsymbol\theta$ includes the intercept and six coefficients.
If the matrix is not full rank, numerical least-squares or a pseudoinverse is
required instead of the simple inverse formula. A negative coefficient is
possible and should be interpreted, not silently forced positive. Neither a
coefficient nor a strong association proves causation.

\section{Completed validation and preliminary model evaluation}
\label{sec:model-evaluation}
\textbf{Data preparation checks already completed:}
\begin{enumerate}
\item \textbf{Consistent season and player keys.} All six factors concern
the 2025--26 regular season. The official GRAV data supply exactly 234
unique player IDs; the merge retains all 234 without duplicate IDs.
\item \textbf{Traded-player aggregation.} The 701-row ON/OFF sources were
combined into player-level factors by ON-minutes-weighted team-stint
differences, with individual factor-specific weights. The check found
24 multi-team players among the official GRAV sample; both resulting
ON/OFF factors are present for all 234 GRAV players.
\item \textbf{Complete-case verification.} Exactly 175 of 234 GRAV players
have all six predictors; the 59 incomplete cases lack qualifying clutch
data under the 30-minute rule. No missing values were imputed.
\end{enumerate}

\textbf{Modeling and evaluation:}
\begin{enumerate}
\item \textbf{Split by player before fitting.} We used an 80\% training /
20\% test split of the 175 complete players: 140 training players and
35 held-out players. The split used NumPy's default random generator with
seed 440. The held-out set was kept separate from model fitting and
cross-validation.
\item \textbf{Keep the test set unseen.} Standardization means and
population standard deviations were computed on training players only
and then applied unchanged to held-out players. Coefficients were fit
without using held-out GRAV scores.
\item \textbf{Fit multiple linear regression by ordinary least squares.}
The six-predictor model was fitted by numerically stable least squares.
For standardized predictors $Z_{ij}$, the fitted equation is
\begin{align}
\widehat Y_i={}&1.5243929543
-0.0232279793Z_{i1}
+0.6215652779Z_{i2}
+0.0428099387Z_{i3}
+1.8882029307Z_{i4}\notag\\
&+2.3815304922Z_{i5}
-0.9446153351Z_{i6}.
\label{eq:fitted-ols}
\end{align}
Predictors $Z_{i1},\ldots,Z_{i6}$ correspond, in order, to
\texttt{OFF\_RATING\_DIFF}, \texttt{EFG\_DIFF\_PP},
\texttt{AST\_PER\_GAME}, \texttt{TIMEOUTS\_PER\_36},
\texttt{CLUTCH\_PTS\_PER\_36}, and \texttt{OWS}. Each predictor is
standardized using its training-set mean and population standard deviation.
\item \textbf{Evaluate prediction accuracy.} For $m$ held-out players,
\begin{equation}
 \mathrm{MAE}=\frac{1}{m}\sum_{i=1}^{m}|Y_i-\widehat Y_i|.
 \label{eq:mae}
\end{equation}
A complementary test-set $R^2$ compares predictions with a reference
mean; held-out $R^2$ can be negative. On this 35-player test set, OLS
achieved MAE $2.8560475861$ GRAV points and $R^2=0.3734245141$.
\item \textbf{Evaluate ranking agreement.} Spearman's rank correlation
$\rho_s$ compares predicted and official GRAV rankings. Without ties,
\begin{equation}
\rho_s=1-\frac{6\sum_{i=1}^{m}d_i^2}{m(m^2-1)},
\label{eq:spearman}
\end{equation}
where $d_i$ is the difference between a player's predicted and official
rank. With ties, compute the Pearson correlation between the two sets
of assigned ranks. The held-out OLS Spearman correlation was
$\rho_s=0.5753501401$.
\item \textbf{Assess training-sample stability.} We performed shuffled
five-fold cross-validation on the 140 training players using seed 440.
Each fold used 112 players for fitting and 28 for validation; predictor
standardization was recomputed from that fold's 112 training players.
Across the five validation folds, the mean and sample standard deviation
were MAE $3.2905\pm0.2153$, test $R^2$ $0.2770\pm0.2887$, and Spearman
$\rho_s$ $0.5300\pm0.2170$. The sizeable fold-to-fold variation,
particularly in $R^2$ and rank correlation, indicates uncertainty in
performance estimates. These fold results use only the 140-player
training set; the 35-player test set remains held out.
\item \textbf{Compare simple baselines.} We compared OLS against the
standardized equal-weight index and single-predictor models. The
equal-weight index was calibrated to GRAV using training players only.
The best single-predictor baseline used clutch points per 36 minutes.
Held-out results were:
\begin{center}
\small
\begin{tabular}{@{}lrrr@{}}
\toprule
\textbf{Model} & \textbf{MAE} & \textbf{Test $R^2$} & \textbf{Spearman $\rho_s$}\\
\midrule
Six-predictor OLS & 2.8560475861 & 0.3734245141 & 0.5753501401\\
Calibrated equal-weight index & 2.8927238477 & 0.3509312110 & 0.4759103641\\
Clutch points/36 only & 2.8589605185 & 0.2921527789 & 0.4893557428\\
\bottomrule
\end{tabular}
\end{center}
OLS was slightly better than the equal-weight benchmark on this particular
split. The 35-player test set is small, and this single split does not
establish that OLS will consistently outperform the benchmark.
\end{enumerate}

\textbf{Reproducibility.} The model is implemented in
\texttt{src/models/train.py} and can be run from the project root with
\texttt{.venv/bin/python src/models/train.py}. The script uses the merged
CSV, seed 440, and a 20\% test fraction by default. Its run also reports
single-predictor baselines. The cross-validation summarizes variation
within the training sample; it does not replace the held-out test
evaluation or establish performance on a new season.

\textbf{Interpretation:} Even accurate held-out predictions would not
reconstruct the NBA's tracking algorithm or show that the six predictor
variables causally create defensive attention.

\subsection{What the evaluation results mean}
The cross-validation value $\rho_s=0.5300$ is Spearman's rank correlation,
not a $p$-value and not a percentage of predictions that were correct. It
indicates a moderately positive average association between the model's
predicted player ordering and the players' official GRAV ordering across
the five validation folds. The standard deviation, $0.2170$, shows that
ranking agreement differed noticeably from fold to fold.

The cross-validation MAE of $3.2905$ means that a validation prediction
was, on average, about 3.29 GRAV points away from the observed score.
Lower MAE is better. The reported $\pm0.2153$ is the sample standard
deviation of the five fold-specific MAEs; it is not a confidence interval.

The mean cross-validation $R^2$ of $0.2770$ indicates that, averaged over
the validation folds, the model reduced squared prediction error by about
27.7\% relative to predicting that fold's mean GRAV for every player.
The fold standard deviation, $0.2887$, is large relative to the mean,
showing substantial variation; individual fold $R^2$ values can be
negative. Thus the average does not imply equally good performance in
every fold.

These cross-validation summaries describe five validation subsets within
the 140-player training sample. Separately, on the untouched 35-player
test set, the model had MAE $2.8560$, $R^2=0.3734$, and
$\rho_s=0.5754$. These are distinct evaluations, not additional folds
to combine with the cross-validation averages. Both suggest some
predictive and ranking signal in this sample, but the small test set and
fold-to-fold variation mean the evidence is preliminary and does not
establish reliable performance on future seasons. The model's associations
also do not show that the predictors cause gravity.

\section{Presentation outline: October 22, 2 p.m.}
\begin{center}
\small
\begin{tabularx}{\linewidth}{@{}cp{0.36\linewidth}Xc@{}}
\toprule
\textbf{Slide} & \textbf{Topic} & \textbf{What to show} & \textbf{Time}\\
\midrule
1 & Gravity and research question & NBA definition; why public proxies are interesting & 1:00\\
2 & Six-factor methodology & Six measured predictors, one deferred, merged CSV & 1:15\\
3 & Mathematical definitions & OffRtg, teammate-only eFG, assists, timeouts, clutch rate, OWS & 2:00\\
4 & Mathematical model & Standardization, OLS objective, equal-weight baseline & 2:00\\
5 & Verified data and model evaluation & 234 merged, 175 complete; held-out and five-fold CV results & 2:00\\
6 & Caveats and next steps & Clutch selection, traded-player proxies, noncausality & 1:15\\
\midrule
 & \textbf{Target total} & & \textbf{9:30}\\
\bottomrule
\end{tabularx}
\end{center}
\textbf{Presentation rule:} Explain that $\rho_s$ measures rank
association, MAE measures average error in GRAV points, and $R^2$
compares squared error with a mean-prediction baseline. Distinguish the
single 35-player held-out results from five-fold cross-validation on the
training players. The fold variation is substantial; do not claim
reliable superiority over the benchmark or imply causal effects.
\textbf{Slide-deck review:} The linked Google Slides deck was inspected in
view-only mode and is not yet aligned with this report. It still presents
open-shot percentage as P1 despite that factor being deferred for lack of
tracking data, and a brainstorming slide proposes Random Forest although
the agreed model is six-predictor OLS. Update these slides before presenting.

\section{Collaboration and remaining responsibilities}
Josue completed the recent acquisition and preparation of the six candidate
factor datasets, the official Gravity target, the traded-player ON/OFF
aggregation, and the final player-level merge. The data audit and this mathematical documentation were prepared to
support the next presentation. An initial six-predictor OLS model has now
been fit and evaluated on one held-out split and by five-fold
cross-validation within the training sample. The results should be
discussed with the group and presented as preliminary. The linked slide
deck was reviewed in view-only mode; an editor with edit access must bring
its hypotheses, variables, and model choice into agreement with this report.

\begin{center}
\small
\begin{tabularx}{\linewidth}{@{}p{0.24\linewidth}Xp{0.19\linewidth}@{}}
\toprule
\textbf{Workstream} & \textbf{Deliverable} & \textbf{Status}\\
\midrule
Data acquisition & Six candidate predictors plus 234 official GRAV scores & Completed\\
Data integration & Aggregated ON/OFF; checked OWS IDs; merged 234 players, 175 complete & Completed\\
Mathematical modeling & Fit OLS with six predictors; evaluate held-out prediction and ranks & Held-out and five-fold CV completed; interpret cautiously\\
Presentation and report & Present verified model results with limitations & Report updated; slide deck needs substantive alignment\\
\bottomrule
\end{tabularx}
\end{center}

\section{Updated decisions and immediate checklist}
\begin{enumerate}
\item[$\boxtimes$] Choose the 2025--26 NBA regular season.
\item[$\boxtimes$] Collect six candidate predictors; defer open-shot data.
\item[$\boxtimes$] Define and calculate opponent timeout responses per 36
season minutes from 1,230 NBA games.
\item[$\boxtimes$] Calculate teammates-only ON/OFF eFG, excluding focal-player
shot attempts from the ON group.
\item[$\boxtimes$] Obtain published \emph{Offensive} Win Shares from
Basketball-Reference and match to NBA player IDs (582 players; 234/234 GRAV matches).
\item[$\boxtimes$] Confirm \texttt{AST} is assists per game; preserve it as
\texttt{AST\_PER\_GAME}.
\item[$\boxtimes$] Extract 234 official numerical GRAV scores from
\texttt{GravityLeaders}, using \texttt{AVGGRAVITYSCORE}.
\item[$\boxtimes$] Aggregate multi-team ON/OFF stints by factor-specific
ON minutes; verify 602 player rows and all 234 GRAV matches.
\item[$\boxtimes$] Check source-minute precision: maximum discrepancy
0.5 minute; use each factor's own source minutes.
\item[$\boxtimes$] Merge all six factors into
\texttt{data/processed/player\_gravity\_model\_data.csv}; verify 234 rows
with 175 six-factor-complete players and 59 lacking clutch data.
\item[$\boxtimes$] Retain P1--P6 in predictor order, with open-shot P7
(TBD) and ranking agreement P8.
\item[$\boxtimes$] Assume all group members can access the merged CSV and
project files through the shared repository, as confirmed by the group.
\item[$\boxtimes$] Fit and evaluate an initial six-predictor OLS model;
record the fitted formula, seed, sample split and held-out results.
\item[$\boxtimes$] Run five-fold cross-validation on the training players,
restandardizing within each fold; report fold means and variability.
\item[$\square$] Revise the linked presentation: defer open-shot P1 until
tracking data exist, remove Random Forest brainstorming, and add the agreed
six-factor OLS method plus held-out and five-fold CV results.
\item[$\square$] Finalize the presentation and its verified findings.
\end{enumerate}

\textbf{Mathematical-report requirement:} All definitions, units,
assumptions, and selection rules must remain explicit. The reported
preliminary fit includes its numerical coefficients and held-out
evaluation measures. Add accurate textbook editions/page references for
least squares, standardization and rank correlation before final submission.

\begin{thebibliography}{9}
\bibitem{nba-gravity-intro}
NBA.com Staff, ``Intro to Gravity,'' updated January 8, 2026. \\
\url{https://www.nba.com/news/intro-to-gravity-stat-nba-2025}

\bibitem{nba-gravity-board}
NBA, ``Inside The Game: Player Gravity Leaders.'' \\
\url{https://www.nba.com/inside-the-game/player/gravity?Season=2025-26&SeasonType=Regular+Season}

\bibitem{nba-gravity-api}
NBA Stats, \emph{GravityLeaders} endpoint, 2025--26 regular-season player
gravity data, retrieved using the Python \texttt{nba\_api} package.
\url{https://stats.nba.com/stats/gravityleaders}

\bibitem{nba-glossary}
NBA, ``Stats Glossary'' (entries for OffRtg and eFG\%). \\
\url{https://www.nba.com/stats/help/glossary}

\bibitem{winshares}
Basketball-Reference, ``NBA Win Shares,'' methodology, especially Section III
(``Crediting Offensive Win Shares to Players''). \\
\url{https://www.basketball-reference.com/about/ws.html}

\bibitem{br-advanced}
Basketball-Reference, ``2025--26 NBA Advanced Stats,'' player-level OWS
column used in the exported and NBA-ID-matched 2025--26 dataset. \\
\url{https://www.basketball-reference.com/leagues/NBA_2026_advanced.html}
\end{thebibliography}

\section*{Shared decision log (updated October 22)}
\begin{center}
\small
\begin{tabularx}{\linewidth}{@{}p{0.29\linewidth}Xp{0.20\linewidth}@{}}
\toprule
\textbf{Question} & \textbf{Current decision / evidence} & \textbf{Status}\\
\midrule
Which factors enter v1? & OffRtg, teammate eFG, assists/game, timeouts/36, clutch points/36, OWS. Open-shot tracking remains deferred. & Six measured\\
What is the target? & Official numerical NBA GRAV (\texttt{AVGGRAVITYSCORE}); 234 distinct players. & Collected\\
Have trades been handled? & Separate ON-minutes-weighted factor 1 and 2 means across team stints, with verified source-minute precision. & Completed\\
Has the dataset been merged? & \texttt{player\_gravity\_model\_data.csv}: 234 rows, 175 fully observed, 59 without qualifying clutch. & Completed locally\\
How are weights chosen? & Six-variable OLS; seed-440 split gave 140 training and 35 held-out players. Held-out OLS MAE 2.8560, $R^2=0.3734$, Spearman 0.5754; five-fold training CV MAE $3.2905\pm0.2153$, $R^2$ $0.2770\pm0.2887$, Spearman $0.5300\pm0.2170$. & Preliminary; fold variability\\
Can teammates read the data? & Group confirms all members have access to the shared GitHub repository and required files; not independently tested here. & Assumed\\
\bottomrule
\end{tabularx}
\end{center}

\textbf{Revision note (October 22, 2026, v0.8).} This version updates
the previous planning draft to reflect completed traded-player aggregation,
minute-discrepancy inspection, confirmed per-game assists, verified
Basketball-Reference OWS provenance and matching, creation of the final
merged 234-player CSV, an initial six-predictor OLS fit, five-fold
cross-validation on the training sample. Held-out and cross-validation
results show fold-to-fold variability and remain preliminary; the linked
presentation remains out of sync: it presents open-shot P1 and Random
Forest brainstorming instead of the agreed OLS model. The project GitHub
repository link is included on the title page.

\end{document}